\section{How the distance bounds are proved}\label{sec:distance-proof}

The minimum-distance event quantifies over every message after setup.
A bound for a fixed message cannot establish this event without an
additional argument. We instead enumerate supports and track the rank
of their zero equations. This section gives the common argument for
Theorems~\ref{thm:width-one-small} and~\ref{thm:wider-guarantees}; its
last subsection states which appendix completes each row of
Tables~\ref{tab:width-one-small} and~\ref{tab:wide-measured}.

\subsection{Branches and equation rank}

For a message $x$, call $B_8x$ its source vector; the support and weight
of $B_8x$ are its \emph{source support} and \emph{source weight}.
Fix a source support $S\subseteq\{0,\ldots,n-1\}$ and define its
shortened source space
\[
 V_S=\{x\in\F_p^k:(B_8x)_i=0\text{ for }i\notin S\},
 \qquad D_S=\dim V_S.
\]
Require every source coordinate in $S$ to be nonzero on the candidate
message. Discard $S$ if $D_S=0$. A \emph{branch} specifies which outputs
are zero in each scan. Fix a forbidden final weight $e$; a branch is
\emph{bad} if its final output has weight at most $e$.

Each permutation assigns to every position of the following scan an
\emph{input label}, the index of the coordinate that position reads.
A label is \emph{active} if its coordinate is required to be nonzero.
We reveal (\emph{expose}) input labels and tap coefficients in scan
order. Every exposed coordinate is a linear form on $V_S$.
Let $W\subseteq V_S^*$ span the zero equations imposed so far, starting
with $W=0$. The same span is retained across both scans.

A required nonzero coordinate must remain a nonzero form in $V_S^*/W$.
Discard a branch whenever this necessary condition fails. Keeping other
infeasible branches only enlarges the count. A zero output whose form
is outside $W$ increases $\dim W$ by one; a zero output whose form lies
in $W$ adds no rank. A \emph{run} is a maximal block of consecutive
nonzero outputs of a scan, and the zero output that ends it is a
\emph{return}. A return whose form already lies in $W$ is a
\emph{free return}.

For $\tau\ge1$, weight a branch of rank $r$ by $\tau^{-r}$.
Any branch realized by a nonzero message has $r\le D_S-1$, because
that message belongs to $\ker W$. Thus, for each fixed setup,
\begin{equation}\label{eq:rank-event-domination}
 \mathbf1_{\{\text{a nonzero message realizes a bad branch from }S\}}
 \le \tau^{D_S-1}\sum_{\gamma\text{ retained and bad}}
                                  \tau^{-r(\gamma)}.
\end{equation}
Several messages may realize the same branch. The right side counts
that branch once, so it introduces no factor depending on $p^{D_S}$.

\subsection{A uniform bound for one scan}

First consider width one. At any point of the exposure, the
\emph{pool} is the multiset of forms at the input labels not yet exposed.
A \emph{continuation bound} is an upper bound on the weighted count of
branch completions from a given feedback state. It must hold for every
quotient $V_S^*/W$, every pool, and every earlier nonzero requirement.
Fresh coefficients remain independent of this exposure history.
Nonzero requirements restrict the branches being counted; they do not
condition the setup distribution.

There are two feedback states: zero (state $0$) and nonzero (state $1$).
Let $x,y>0$ be formal markers for nonzero inputs and outputs; from here
on, $x$ and $y$ denote these markers, not a message or a scan output.
Let $v_0,v_1>0$ be continuation bounds for the two states.
The following positive homogeneous operator bounds the two states:
\begin{align}
 T_0(v)&=v_0+xyv_1,\nonumber\\
 T_1(v)&=yv_1+x\max\left\{
        yv_1+\frac{v_0}{\tau},\;
        \left(1-\frac1q\right)yv_1+\frac{v_0}{q}
                         \right\}.\label{eq:width-one-quotient}
\end{align}
To justify the maximum, consider a nonzero input form and a nonzero
feedback form in the quotient by $W$, before exposing the fresh tap.
If these forms are independent, requiring a zero output adds one
independent equation. This gives the first alternative. If the forms
are dependent, at most one coefficient produces the zero form.
The zero and nonzero outcomes are then mutually exclusive, giving the
second alternative when this coefficient belongs to the alphabet.
If it does not, the first alternative bounds the nonzero outcome.
For $q=1$ the only coefficient is one, so a dependent zero output is
deterministic, the second alternative is $v_0$, and only the two
permutations are random.

Normalize $v_0=1$, set $\lambda=\max_i T_i(v)/v_i$, and
$C=\max_i1/v_i$. Monotonicity and homogeneity imply
$[T^n(\mathbf1)]_0\le C\lambda^n$. For sparse supports we instead use
\begin{equation}\label{eq:scan-terminal-sixteen}
 [T^n(\mathbf1)]_0\le
 \lambda^{n-16}\max_i\frac{[T^{16}(\mathbf1)]_i}{v_i}.
\end{equation}
The terminal iterations are evaluated exactly over the rationals.

\subsection{Composing the scans}

For the Newton precode, put
\begin{equation}\label{eq:newton-source-polynomial}
 Q_8(z,\tau)=1+\sum_{j=9}^{16}\binom{16}{j}z^j\tau^{j-8},
 \qquad F_s(\tau)=[z^s]Q_8(z,\tau)^{n/16}.
\end{equation}
The coefficient $F_s(\tau)$ sums $\tau^{D_S}$ over source supports of
size $s$. Indeed, a support of size $j\ge9$ within one block has
shortening dimension $j-8$ by the MDS property, and there are
$\binom{16}{j}$ such supports. The constant term represents a zero block.

Fix the incoming forms of a scan, and suppose $\kappa$ of its input
labels are active. Give a branch with output weight $h'$ and incremental
rank $r'$ the weight $x^\kappa y^{h'}\tau^{-r'}$. Average this weight
over the order of the distinct active labels and over all fresh
coefficients, then sum over placements of the active labels and over
output-support branches. A \emph{uniform scan bound} is a function
$K(\kappa)$ that bounds this sum for every incoming quotient and every
set of earlier nonzero requirements. The two-state bound above gives
$K(\kappa)=C\lambda^n$. The bounds for plain accumulators also average
over the active labels and may include a $\kappa$-dependent factor.

\begin{lemma}[Composition of uniform scan bounds]\label{lem:scan-composition}
Fix source weight $s$, intermediate weight $h$, and forbidden final
weight $e$. For scan $j$, let $K_j$ be a uniform scan bound, with
markers $x_j,y_j>0$, $y_2\le1$, and the same $\tau\ge1$ in both scans.
The probability of a nonzero message with these source and intermediate
weights and final weight at most $e$ is at most
\begin{equation}\label{eq:general-profile}
 \frac{F_s(\tau)}{\tau}
 \frac{K_1(s)K_2(h)}
 {\binom ns\binom nh\,x_1^s(y_1x_2)^h y_2^e}.
\end{equation}
\end{lemma}

\begin{proof}
Fix $S$ of size $s$. A uniform permutation sends its active labels to
a uniform $s$-subset and gives them a uniform order on that subset.
These choices are independent. The first normalization is therefore
$1/\binom ns$; the label-order average is already included in $K_1(s)$.

Fix a complete first-scan history with intermediate weight $h$. The
second permutation and its coefficients remain independent of that
history. Apply $K_2(h)$ in the quotient left by the first scan.
The second support-placement normalization is $1/\binom nh$.
Rank increments add, so the product of the two incremental weights is
the single weight $\tau^{-r}$ in~\eqref{eq:rank-event-domination}.
This conditional argument does not assume independent zero equations.

All terms are nonnegative. Extracting source and intermediate weights
divides by $x_1^s(y_1x_2)^h$. For a final weight $j\le e$,
$y_2^j\ge y_2^e$, which gives the remaining divisor. Apply
inequality~\eqref{eq:rank-event-domination}, then sum over $S$; this
supplies $F_s(\tau)/\tau$.
\end{proof}

For $K_j=C_j\lambda_j^n$, the lemma gives the width-one profile bound
\begin{equation}\label{eq:width-one-profile}
 \frac{F_s(\tau)}{\tau}
 \frac{C_1\lambda_1^n C_2\lambda_2^n}
 {\binom ns\binom nh\,x_1^s(y_1x_2)^h y_2^e}.
\end{equation}
Uniformity is essential here: a bound proved only for a fresh,
unconstrained input space would not justify the second application.

\subsection{Finite evaluation and the plain-accumulator refinements}

The inverse width-one scan gives $h\le\min(n,2e)$ and
$s\le\min(n,2h)$. The \emph{direct bound} evaluates~\eqref{eq:width-one-profile}
with the operator~\eqref{eq:width-one-quotient} on every such integer
profile. For source weights below 32, it uses the exact coefficients
$F_s(\tau)$ and either the terminal bound~\eqref{eq:scan-terminal-sixteen}
or the hypergeometric support bound of Appendix~\ref{app:hyper-support}.
The remaining weights use $F_s\le Q_8(z,\tau)^{n/16}/z^s$ and convex
bounds on rectangles in the $(s,h)$ plane. These rectangles come from a
recursive subdivision of the profile domain; each final rectangle is a
\emph{leaf}. Each leaf is charged its bound times its number of integer
profiles. The parameters for each leaf are chosen numerically, and every
resulting inequality is checked in rational arithmetic with directed
logarithm bounds. The check confirms that the leaves cover the domain and
sums their upward-rounded contributions. The rows of
Table~\ref{tab:width-one-small} also use the stronger bounds described next.

The main difficulty at $q=1$ is the dependent case of~\eqref{eq:width-one-quotient}.
For some pools, a single return is free with probability one.
Maximizing that probability separately at every step ignores the cost
of forming such pools. The refinements for plain accumulators instead
charge several exposures jointly.

For example, draw two nonzero forms independently from one pool.
Let $P_\parallel$ be the probability that they are proportional and
$P_-$ the probability that they are exact opposites. Then
$P_-\le P_\parallel/2$: in odd characteristic the nonzero forms on each
line split into disjoint pairs $\{f,-f\}$, so a line of total mass
$\sigma_\ell$ yields an opposite pair with probability at most
$\sigma_\ell^2/2$, while $P_\parallel=\sum_\ell\sigma_\ell^2$.
A proportional nonopposite pair cannot sum to zero on a surviving
branch: imposing that equation would kill both required inputs.
The \emph{projective-pair bound} (Appendix~\ref{app:fresh-pairs}) uses
this restriction for the first two inputs of each run. The
\emph{joint-renewal bound} (Appendix~\ref{app:joint-renewal}) retains it
through neighboring returns. The \emph{complete-run bound}
(Appendix~\ref{app:complete-runs}) counts a whole run, with its free
returns and rank-gaining returns, in a fixed quotient, then lets the
quotient change when a new zero equation is imposed. Each bound is
uniform over every pool, so Lemma~\ref{lem:scan-composition} applies
after each scan.

Independent draws are an auxiliary counting device. The actual
permutation draws labels without replacement. The complete-run bound
absorbs the resulting correction in an explicit factor $\kappa^{16}$
for a scan with $\kappa$ active inputs. Thus its two-scan bound uses
$K_1(s)K_2(h)=(sh)^{16}C_1\lambda_1^nC_2\lambda_2^n$. It also includes
the correction terms that arise because $p$ is finite and the runs still
open at the end of a scan.

\subsection{Which appendix proves each row}

The $q=1$ row of Table~\ref{tab:width-one-small} is proved in
Appendix~\ref{app:complete-runs}. Its cover assigns each leaf the
projective-pair bound of Appendix~\ref{app:fresh-pairs}, the
joint-renewal bound of Appendix~\ref{app:joint-renewal}, or the
complete-run bound, and it bounds source weights below 32 by the
hypergeometric support bound of Appendix~\ref{app:hyper-support}.
The rows with $q\ge2$ are proved in Appendix~\ref{app:small-q-pairs},
which extends projective pairs to small coefficients and chooses, on
each leaf, either that bound or the direct bound. For source weights
below 32 it uses the hypergeometric support bound at $q=2$ and the
direct bound with~\eqref{eq:scan-terminal-sixteen} for $q\ge4$.
The two rows of Table~\ref{tab:wide-measured} are proved in
Appendix~\ref{app:wide-mask}, which replaces~\eqref{eq:width-one-quotient}
by an operator on support masks and ranks of the preceding $w$ outputs
and treats source weights below 32 with exact source coefficients and
sixteen terminal iterations. Each certificate covers every sparse weight
and dense support profile.
